\begin{theorem}[Alternating series test (Leibniz criterion)] \label{theorem:alternating_series_test_for_decreasing_magnitude_terms}
    \TODO{AI generated, verify}
    Let $(a_n)_{n=1}^\infty$ be a sequence in $\mathbb{R}$\CrefIfExists{definition:field_of_real_numbers} such that:
    \begin{enumerate}
        \item $a_n \geq 0$ for all $n$;
        \item The sequence is monotonically nonincreasing, i.e., $a_{n+1} \leq a_n$ for all $n$;
        \item $\lim_{n \to \infty} a_n = 0$.
    \end{enumerate}
    Then the alternating series $\sum_{n=1}^\infty (-1)^{n-1} a_n$ converges.

    Moreover, for every $N \in \mathbb{N}$, the remainder satisfies
    $$\left| \sum_{n=N+1}^\infty (-1)^{n-1} a_n \right| \leq a_{N+1}.$$
\end{theorem}
\begin{proof}
    \TODO{AI generated, verify}
    Let $S_N = \sum_{n=1}^N (-1)^{n-1} a_n$ denote the partial sums. We show that $(S_N)$ converges by proving that the subsequences of even and odd partial sums converge to the same limit.

    Consider consecutive partial sums:
    $$S_{2m} - S_{2m-2} = (-1)^{2m-2} a_{2m-1} + (-1)^{2m-1} a_{2m} = a_{2m-1} - a_{2m}.$$

    Since $a_n$ is monotonically nonincreasing, we have $a_{2m-1} \geq a_{2m}$, so $S_{2m} \geq S_{2m-2}$. Thus $(S_{2m})_m$ is monotonically nondecreasing.

    Similarly:
    $$S_{2m+1} - S_{2m-1} = (-1)^{2m-1} a_{2m} + (-1)^{2m} a_{{2m+1}} = -a_{2m} + a_{2m+1}.$$

    Since $a_{2m} \geq a_{2m+1}$, we have $S_{2m+1} \leq S_{2m-1}$. Thus $(S_{2m+1})_m$ is monotonically nonincreasing.

    We now establish boundedness. For any $m$, consider the relationship between even and odd partial sums:
    $$S_{2m+1} - S_{2m} = (-1)^{2m} a_{2m+1} = a_{2m+1} \geq 0.$$

    Therefore $S_{2m+1} \geq S_{2m}$ for all $m$. Since $(S_{2m})$ is nondecreasing and $(S_{2k+1})$ is nonincreasing, we have for any $j, k$:
    $$S_2 \leq S_{2j} \leq S_{2k+1} \leq S_1.$$

    The even partial sums are bounded above by $S_1 = a_1$, and the odd partial sums are bounded below by $S_2 = a_1 - a_2$. By the monotone convergence theorem, both subsequences converge:
    $$\lim_{m \to \infty} S_{2m} = L_{\text{even}}, \qquad \lim_{m \to \infty} S_{2m+1} = L_{\text{odd}}.$$

    The difference between consecutive partial sums is:
    $$S_{2m+1} - S_{2m} = a_{2m+1}.$$

    Taking the limit as $m \to \infty$ and using $\lim_{n \to \infty} a_n = 0$:
    $$L_{\text{odd}} - L_{\text{even}} = \lim_{m \to \infty} a_{2m+1} = 0.$$

    Hence $L_{\text{odd}} = L_{\text{even}}$, so $(S_N)$ converges to this common limit.

    For the remainder estimate, let $S = \lim_{N \to \infty} S_N$. We have:
    $$S - S_N = \sum_{n=N+1}^\infty (-1)^{n-1} a_n.$$

    If $N$ is even (say $N = 2m$), then the tail begins with a positive term:
    $$S - S_{2m} = a_{2m+1} - a_{2m+2} + a_{2m+3} - \cdots.$$

    Grouping terms as $(a_{2m+1} - a_{2m+2}) + (a_{2m+3} - a_{2m+4}) + \cdots$, each grouped pair is nonnegative, so $S - S_{2m} \geq 0$. On the other hand, writing as $a_{2m+1} - (a_{2m+2} - a_{2m+3}) - (a_{2m+4} - a_{2m+5}) - \cdots$, each subtracted pair is nonnegative, so $S - S_{2m} \leq a_{2m+1}$. Therefore:
    $$0 \leq S - S_{2m} \leq a_{2m+1}.$$

    If $N$ is odd (say $N = 2m-1$), the tail begins with a negative term. By an analogous argument, we obtain:
    $$-a_{2m} \leq S - S_{2m-1} \leq 0.$$

    In both cases, $|S - S_N| \leq a_{N+1}$, which is the desired remainder estimate.
\end{proof}
